\documentclass[11pt]{article}
\usepackage[margin=1in]{geometry}
\usepackage{amsmath,amssymb,mathtools,amsthm}
\usepackage{booktabs,longtable,array}
\usepackage{enumitem}
\usepackage[hidelinks]{hyperref}

\newcommand{\R}{\mathbb{R}}
\newcommand{\dd}{\mathrm{d}}
\newcommand{\TT}{\mathrm{TT}}
\newcommand{\ones}{\mathbf 1}
\newcommand{\diag}{\operatorname{diag}}
\newcommand{\radinv}{\operatorname{radinv}}
\newcommand{\vecop}{\operatorname{vec}}

\newtheorem{proposition}{Proposition}
\newtheorem{lemma}{Lemma}
\theoremstyle{definition}
\newtheorem{definition}{Definition}
\newtheorem{assumption}{Assumption}
\theoremstyle{remark}
\newtheorem{remark}{Remark}

% The specification cross-references many displayed formulas.  Number all
% display math so references point to visible equations rather than section
% numbers.
\renewcommand{\[}{\begin{equation}}
\renewcommand{\]}{\end{equation}}

\title{An Implementable Fixed-Branch Squared Tensor-Train Filter}
\author{BayesFilter P15 Zhao--Cui TT Specification}
\date{2026-05-31}

\begin{document}
\maketitle

\begin{abstract}
This note specifies one concrete fixed-branch squared tensor-train filtering
algorithm and its analytical same-scalar gradient.  It is inspired by the
sequential squared tensor-train construction of Zhao and Cui, but it is not a
description of every adaptive choice in their research code.  The purpose is
more modest and more operational: define a single value path that a reader can
implement without opening any other source.  The fixed path uses affine domain
maps, normalized Legendre bases, deterministic Halton fitting points, fixed
ranks, ridge alternating least squares, a uniform defensive density, a frozen
shift, explicit squared-TT marginalization, and a finite-difference parity
test.  The resulting scalar is an approximate likelihood; the gradient derived
here differentiates that approximate scalar on the saved branch.
\end{abstract}

\tableofcontents

\section{What This Specification Implements}

The algorithm in this note takes a state-space model and produces an
approximate filtering likelihood.  The model has state \(x_t\), optional static
parameter \(\vartheta\), observation \(y_t\), transition density
\(f_\alpha(x_t\mid x_{t-1},\vartheta)\), and observation density
\(g_\alpha(y_t\mid x_t,\vartheta)\).  The external parameter \(\alpha\) is the
quantity with respect to which the approximate likelihood will be
differentiated.

At each time \(t\), the method constructs an approximation to the unnormalized
filtering density over the augmented variable
\[
  r_t=(x_t,\vartheta,x_{t-1})\in\R^D.
\]
The approximation has the form
\[
  \widehat q_t(z;\alpha;B_t)
  =
  \phi_t(z;\alpha;B_t)^2+\tau_t 2^{-D},
  \qquad z\in[-1,1]^D,
  \label{eq:p15-qhat}
\]
where \(B_t\) denotes the fixed branch data.  The square-root tensor train
\(\phi_t\) is fitted by a fully specified fixed-rank ridge ALS procedure.  The
defensive term \(\tau_t2^{-D}\) is a small uniform density floor on the
reference box.

The one-step approximate evidence contribution is
\[
  \widehat z_t(\alpha;B_t)
  =
  \int_{[-1,1]^D}\widehat q_t(z;\alpha;B_t)\,\dd z
  =
  R_{t,0}(\alpha;B_t)+\tau_t,
  \label{eq:p15-zhat}
\]
and the declared approximate log-likelihood is
\[
  \widehat\ell_T^{\TT}(\alpha;B)
  =
  \sum_{t=1}^T
  \{\log(R_{t,0}(\alpha;B_t)+\tau_t)-c_t\}.
  \label{eq:p15-declared-scalar}
\]
The constants \(c_t\) are frozen stabilizing shifts stored in the branch.  The
gradient derived below is the derivative of \eqref{eq:p15-declared-scalar} for
the saved branch \(B=(B_1,\ldots,B_T)\).  It is not the derivative of an
adaptive algorithm whose ranks, fitting points, shifts, domains, or
truncations change with \(\alpha\).

This distinction matters.  A filtering algorithm may adapt its numerical
representation as the parameter changes.  That can be useful for accuracy, but
then the formula being evaluated changes.  A classical analytical gradient
must differentiate one declared formula.  P15 therefore freezes the branch
objects and differentiates the resulting approximate scalar.

\section{State-Space Filtering Object And Approximate Likelihood}

Assume that before observing \(y_t\), the previous approximate filter is a
normalized density \(\widehat p_{t-1}(x_{t-1},\vartheta;\alpha)\).  The exact
Bayesian one-step unnormalized joint density would be
\[
  q_t(r_t;\alpha)
  =
  \widehat p_{t-1}(x_{t-1},\vartheta;\alpha)
  f_\alpha(x_t\mid x_{t-1},\vartheta)
  g_\alpha(y_t\mid x_t,\vartheta).
  \label{eq:p15-qt}
\]
Its integral is
\[
  Z_t(\alpha)
  =
  \int q_t(x_t,\vartheta,x_{t-1};\alpha)
  \,\dd x_t\,\dd\vartheta\,\dd x_{t-1}.
  \label{eq:p15-exact-Z}
\]
Integrating first over \(x_{t-1}\) gives the prediction density, and then
integrating against the likelihood gives the innovation evidence:
\[
\begin{aligned}
  Z_t(\alpha)
  &=
  \int g_\alpha(y_t\mid x_t,\vartheta)
  \left[
  \int f_\alpha(x_t\mid x_{t-1},\vartheta)
       \widehat p_{t-1}(x_{t-1},\vartheta;\alpha)\,\dd x_{t-1}
  \right]\dd x_t\,\dd\vartheta .
\end{aligned}
  \label{eq:p15-filtering-integral}
\]
The updated filter is the marginal
\[
  p_t(x_t,\vartheta;\alpha)
  =
  \int
  \frac{q_t(x_t,\vartheta,x_{t-1};\alpha)}
       {Z_t(\alpha)}
  \,\dd x_{t-1}.
  \label{eq:p15-exact-filter-marginal}
\]

The P15 method replaces \(q_t\) by a nonnegative finite-dimensional
approximation \(\widehat q_t\), integrates \(\widehat q_t\) to get
\(\widehat z_t\), and marginalizes \(\widehat q_t/\widehat z_t\) to get the
next approximate filter.  Thus the connection to filtering is not decorative:
the tensor train is only a way to compute the two integrals in
\eqref{eq:p15-filtering-integral} and \eqref{eq:p15-exact-filter-marginal}.

\section{Fixed Computational Domain And Reference Measure}

Choose a hyperrectangle
\[
  \Omega_t=\prod_{k=1}^D[\ell_{t,k},u_{t,k}]
\]
with finite endpoints.  The reference coordinate is \(z\in[-1,1]^D\).  Define
\[
  a_{t,k}=\frac{u_{t,k}+\ell_{t,k}}2,\qquad
  h_{t,k}=\frac{u_{t,k}-\ell_{t,k}}2,\qquad
  r_{t,k}=a_{t,k}+h_{t,k}z_k.
  \label{eq:p15-affine-map}
\]
Let \(\Psi_t(z)\) denote this affine map.  The Jacobian determinant is
\[
  J_t=\prod_{k=1}^D h_{t,k}.
  \label{eq:p15-jacobian}
\]
The transformed unnormalized density is
\[
  \widetilde q_t(z;\alpha)=q_t(\Psi_t(z);\alpha)J_t.
  \label{eq:p15-transformed-q}
\]

\begin{proposition}[Affine-domain transformation]
For any integrable \(q_t\) on \(\Omega_t\),
\[
  \int_{\Omega_t}q_t(r;\alpha)\,\dd r
  =
  \int_{[-1,1]^D}\widetilde q_t(z;\alpha)\,\dd z.
  \label{eq:p15-change-vars}
\]
\end{proposition}

\begin{proof}
Equation \eqref{eq:p15-affine-map} is diagonal affine.  Hence
\(\dd r=J_t\,\dd z\).  Substituting \(r=\Psi_t(z)\) gives
\[
  \int_{\Omega_t}q_t(r;\alpha)\,\dd r
  =
  \int_{[-1,1]^D}q_t(\Psi_t(z);\alpha)J_t\,\dd z,
\]
which is \eqref{eq:p15-change-vars}.
\end{proof}

For the fixed-branch derivative, \(\ell_{t,k},u_{t,k},a_{t,k},h_{t,k}\), and
\(J_t\) are frozen.  If a future implementation adapts the box with
\(\alpha\), that is a different scalar and requires additional boundary and
Jacobian derivatives.

\section{Legendre Basis, Evaluation Recurrence, And Mass Matrices}

On each coordinate, use normalized Legendre functions
\[
  b_{k,a}(z)=\sqrt{\frac{2a+1}{2}}\,P_a(z),
  \qquad a=0,\ldots,p_k-1,
  \label{eq:p15-legendre-basis}
\]
where \(P_a\) is the degree-\(a\) Legendre polynomial.  The implementation
evaluates them by
\[
  P_0(z)=1,\qquad P_1(z)=z,
  \label{eq:p15-legendre-start}
\]
\[
  (a+1)P_{a+1}(z)=(2a+1)zP_a(z)-aP_{a-1}(z),
  \qquad a\ge1.
  \label{eq:p15-legendre-recurrence}
\]
For a vector \(z_k^{(1)},\ldots,z_k^{(N)}\), store the basis matrix
\[
  B_k[j,a]=b_{k,a}(z_k^{(j)}).
  \label{eq:p15-basis-matrix}
\]

\begin{proposition}[Legendre mass matrix]
The one-dimensional mass matrix
\[
  M_k[a,b]=\int_{-1}^{1}b_{k,a}(z)b_{k,b}(z)\,\dd z
  \label{eq:p15-mass}
\]
equals the identity matrix.
\end{proposition}

\begin{proof}
Legendre orthogonality gives
\[
  \int_{-1}^{1}P_a(z)P_b(z)\,\dd z
  =
  \frac{2}{2a+1}\mathbf 1_{\{a=b\}}.
\]
Multiplying by the normalization constants in \eqref{eq:p15-legendre-basis}
gives
\[
  M_k[a,b]
  =
  \sqrt{\frac{2a+1}{2}}\sqrt{\frac{2b+1}{2}}
  \frac{2}{2a+1}\mathbf 1_{\{a=b\}}
  =
  \mathbf 1_{\{a=b\}}.
\]
\end{proof}

This identity is why the P15 reference path uses normalized Legendre bases.
Mass contractions then become ordinary Euclidean coefficient contractions.

\section{Tensor-Train Functions And Squared-TT Densities}

A scalar function on the reference box is represented by cores
\[
  C_{k,a}\in\R^{r_{k-1}\times r_k},\qquad r_0=r_D=1,
  \label{eq:p15-core-coefficients}
\]
and one-dimensional matrix-valued factors
\[
  G_k(z_k)=\sum_{a=0}^{p_k-1}C_{k,a}b_{k,a}(z_k).
  \label{eq:p15-G-factor}
\]
The tensor-train function is
\[
  \phi(z)=G_1(z_1)G_2(z_2)\cdots G_D(z_D),
  \label{eq:p15-tt-function}
\]
where the product is a \(1\times1\) scalar.

The density approximation uses the square
\[
  \widehat q_t(z)=\phi_t(z)^2+\tau_t2^{-D}.
  \label{eq:p15-square-density}
\]
The square guarantees nonnegativity.  The uniform term prevents a zero density
floor on the compact reference box.  It does not make a bad fit accurate; it
only ensures that the approximate density remains a proper nonnegative object
after normalization.

\section{Fixed Fitting Design Points}

The fitting design is deterministic.  Let \(p_1,\ldots,p_D\) be the first
\(D\) primes.  If
\[
  j=\sum_{m=0}^{M(j,p)}d_m p^m,\qquad d_m\in\{0,\ldots,p-1\},
\]
define the radical inverse
\[
  \radinv_p(j)=\sum_{m=0}^{M(j,p)}d_m p^{-(m+1)}.
  \label{eq:p15-radinv}
\]
For \(j=1,\ldots,N_{\rm fit}\), the fitting point is
\[
  z_k^{(j)}=2\radinv_{p_k}(j)-1,\qquad k=1,\ldots,D.
  \label{eq:p15-halton}
\]
The least-squares weight matrix is
\[
  W=\frac1{N_{\rm fit}}I.
  \label{eq:p15-uniform-W}
\]
The points, primes, and weights are saved branch objects.  They are not
resampled when evaluating \(\widehat\ell_T^{\TT}(\alpha\pm h e_i;B)\) in a
finite-difference check.

\section{Fixed-Rank Ridge ALS Construction Of The Square-Root TT}

The target values at the fixed design points are
\[
  y_j(\alpha)
  =
  \exp(c_t/2)\sqrt{\widetilde q_t(z^{(j)};\alpha)},
  \qquad j=1,\ldots,N_{\rm fit}.
  \label{eq:p15-target-y}
\]
The stabilizing shift \(c_t\) is frozen and defined at a nominal value
\(\alpha_0\) by
\[
  c_t=-\max_{1\le j\le N_{\rm fit}}\log \widetilde q_t(z^{(j)};\alpha_0).
  \label{eq:p15-shift}
\]
The target values therefore remain numerically moderate near the branch point.
The shift enters the declared scalar as \(-c_t\), as in
\eqref{eq:p15-declared-scalar}.

Fix ranks \(r_0=r_D=1\) and \(r_1,\ldots,r_{D-1}\).  Fix ridge
\(\rho>0\), number of bidirectional ALS sweeps \(S\), and the following exact
initialization.  Let \(e^{(m)}_1\) be the first coordinate vector in
\(\R^m\).  The constant slice of core \(k\) is
\[
  C_{k,0}^{(0)}
  =
  \gamma_k\, e^{(r_{k-1})}_1 (e^{(r_k)}_1)^\top,
  \qquad
  C_{k,a}^{(0)}=0\quad(a\ge1),
  \label{eq:p15-initialization}
\]
where
\[
  \gamma_1=\bar y,\qquad
  \gamma_k=1\quad(k=2,\ldots,D),\qquad
  \bar y=\frac1{N_{\rm fit}}\sum_{j=1}^{N_{\rm fit}}y_j.
  \label{eq:p15-ybar}
\]
Thus only the first rank channel is active initially.  All other rank channels
start at zero and can become nonzero only through the declared ridge ALS
updates.  This convention fixes the branch even when ranks exceed one.

During an ALS update of core \(k\), all other cores are fixed.  Define the left
environment
\[
  L_j^{(k)}=G_1(z_1^{(j)})\cdots G_{k-1}(z_{k-1}^{(j)})
  \in\R^{1\times r_{k-1}}
  \label{eq:p15-left-env}
\]
and the right environment
\[
  R_j^{(k)}=G_{k+1}(z_{k+1}^{(j)})\cdots G_D(z_D^{(j)})
  \in\R^{r_k\times1}.
  \label{eq:p15-right-env}
\]
If \(g_k=\vecop(C_{k,0},\ldots,C_{k,p_k-1})\) is ordered
lexicographically in \((\ell_{\rm left},a,\ell_{\rm right})\), with the right
rank varying fastest, then the row that maps \(g_k\) to \(\phi(z^{(j)})\) is
\[
  A_{j,\cdot}^{(k)}
  =
  (R_j^{(k)})^\top\otimes
  b_k(z_k^{(j)})^\top\otimes
  L_j^{(k)}.
  \label{eq:p15-als-row}
\]
The fixed ridge update is
\[
  N_k g_k=d_k,\qquad
  N_k=(A^{(k)})^\top W A^{(k)}+\rho I,\qquad
  d_k=(A^{(k)})^\top W y.
  \label{eq:p15-ridge-normal}
\]
The algorithm applies updates for \(k=1,\ldots,D\) and then
\(k=D,\ldots,1\), repeated for exactly \(S\) sweeps.  There is no rank
adaptation, no pivot selection, and no truncation in the declared branch.

\section{Normalizer, Defensive Density, Shift, And Approximate Log-Likelihood}

The defensive density is uniform:
\[
  \lambda(z)=2^{-D},\qquad
  \int_{[-1,1]^D}\lambda(z)\,\dd z=1.
  \label{eq:p15-uniform-lambda}
\]
Choose an \(\epsilon_\tau>0\) and set
\[
  \tau_t=\epsilon_\tau 2^D.
  \label{eq:p15-tau}
\]
This makes the defensive floor equal to \(\epsilon_\tau\):
\[
  \tau_t 2^{-D}=\epsilon_\tau.
  \label{eq:p15-floor}
\]
For the primary P15 scalar, \(\tau_t\) and \(c_t\) are fixed branch constants,
so
\[
  \partial_i\tau_t=0,\qquad \partial_i c_t=0.
  \label{eq:p15-tau-c-derivatives}
\]

To integrate \(\phi_t^2\), start with \(R_D=1\) and compute backward
\[
  R_{k-1}
  =
  \sum_{a,b=0}^{p_k-1}
  M_k[a,b]\,
  C_{k,a}R_kC_{k,b}^{\top}.
  \label{eq:p15-R-recursion}
\]
Because \(M_k=I\), this becomes
\[
  R_{k-1}
  =
  \sum_{a=0}^{p_k-1}
  C_{k,a}R_kC_{k,a}^{\top}.
  \label{eq:p15-R-legendre}
\]

\begin{proposition}[Squared-TT mass contraction]
The recursion \eqref{eq:p15-R-recursion} satisfies
\[
  R_0=\int_{[-1,1]^D}\phi_t(z)^2\,\dd z.
  \label{eq:p15-R0-integral}
\]
\end{proposition}

\begin{proof}
For the last coordinate,
\[
  \int G_D(z_D)G_D(z_D)^\top\,\dd z_D
  =
  \sum_{a,b}C_{D,a}C_{D,b}^\top
  \int b_{D,a}(z_D)b_{D,b}(z_D)\,\dd z_D,
\]
which is \(R_{D-1}\).  Suppose \(R_k\) equals the suffix integral over
coordinates \(k+1,\ldots,D\).  Multiplying by \(G_k(z_k)\) on the left and
\(G_k(z_k)^\top\) on the right and integrating \(z_k\) gives exactly
\eqref{eq:p15-R-recursion}.  Induction to \(k=1\) gives
\eqref{eq:p15-R0-integral}.
\end{proof}

Thus
\[
  \widehat z_t=R_0+\tau_t,\qquad
  \widehat\ell_t^{\TT}=\log(R_0+\tau_t)-c_t.
  \label{eq:p15-step-ell}
\]

\section{Marginalization In Squared-TT Form And The Next Filter Object}

Let \(K\) be the set of kept coordinates, such as \(x_t,\vartheta\), and
\(O\) the set of old-state coordinates to integrate out.  Define square cores
\[
  H_k(z_k)=G_k(z_k)\otimes G_k(z_k).
  \label{eq:p15-square-core}
\]
Their coefficient slices are
\[
  H_{k,(a,b)}=C_{k,a}\otimes C_{k,b}.
  \label{eq:p15-square-core-coeff}
\]
For an integrated coordinate \(k\in O\), replace \(H_k\) by
\[
  \mathcal I_k
  =
  \int_{-1}^{1}H_k(z_k)\,\dd z_k
  =
  \sum_{a,b}M_k[a,b]\,C_{k,a}\otimes C_{k,b}.
  \label{eq:p15-integrated-square-core}
\]
For a kept coordinate \(k\in K\), evaluate \(H_k(z_k)\).  Contract the chain
in the original order, using \(\mathcal I_k\) for integrated coordinates.  The
square-term marginal numerator is
\[
  a_t^{\rm sq}(z_K)
  =
  \int \phi_t(z)^2\,\dd z_O.
  \label{eq:p15-square-marginal}
\]
The defensive contribution is
\[
  a_t^{\rm def}(z_K)
  =
  \tau_t 2^{-D}\int_{[-1,1]^{|O|}}1\,\dd z_O
  =
  \tau_t2^{-|K|}.
  \label{eq:p15-defensive-marginal}
\]
The stored numerator and filter are
\[
  a_t(z_K)=a_t^{\rm sq}(z_K)+\tau_t2^{-|K|},
  \qquad
  \widehat p_t(z_K)=a_t(z_K)/\widehat z_t.
  \label{eq:p15-filter-object}
\]

\begin{proposition}[Squared-TT marginalization]
Equation \eqref{eq:p15-filter-object} is the marginal of
\(\widehat q_t/\widehat z_t\) over the old coordinates.
\end{proposition}

\begin{proof}
For the square part, expand \(\phi_t^2\) as the product of square cores
\eqref{eq:p15-square-core}.  Fubini's theorem allows each integrated
coordinate \(k\in O\) to be replaced by \(\mathcal I_k\) in
\eqref{eq:p15-integrated-square-core}.  Kept coordinates are not integrated
and therefore remain as functions of \(z_K\).  This gives
\eqref{eq:p15-square-marginal}.  The uniform defensive density is constant, so
integrating \(|O|\) coordinates multiplies it by \(2^{|O|}\), giving
\eqref{eq:p15-defensive-marginal}.  Dividing by \(\widehat z_t\) normalizes the
marginal.
\end{proof}

The next-step fitting routine needs pointwise values.  The saved filter object
therefore stores a concrete contraction scheme.  First build the square-core
coefficient slices
\[
  \mathcal H_{k,ab}=C_{k,a}\otimes C_{k,b}
  \in\R^{(r_{k-1}^2)\times(r_k^2)}.
  \label{eq:p15-H-coeff-shape}
\]
For an integrated coordinate define the constant matrix
\[
  Q_k=\sum_{a,b}M_k[a,b]\mathcal H_{k,ab}.
  \label{eq:p15-Qk}
\]
For a kept coordinate and a point \(z_k\), define
\[
  E_k(z_k)=\sum_{a,b}b_{k,a}(z_k)b_{k,b}(z_k)\mathcal H_{k,ab}.
  \label{eq:p15-Ek}
\]
Now set
\[
  T_k(z_K)=
  \begin{cases}
    E_k(z_k), & k\in K,\\
    Q_k, & k\in O.
  \end{cases}
  \label{eq:p15-Tk}
\]
The square marginal is the scalar chain contraction
\[
  a_t^{\rm sq}(z_K)=T_1(z_K)T_2(z_K)\cdots T_D(z_K).
  \label{eq:p15-a-sq-chain}
\]
The object also stores prefix and suffix contractions
\[
  P_0=1,\quad P_k=P_{k-1}T_k,
  \qquad
  S_D=1,\quad S_{k-1}=T_kS_k,
  \label{eq:p15-prefix-suffix}
\]
so that an implementation can evaluate the numerator in one left-to-right
pass and, when needed, reuse environments for derivatives.  The value and
log-value evaluators are
\[
  a_t(z_K)=a_t^{\rm sq}(z_K)+\tau_t2^{-|K|},
  \qquad
  \log\widehat p_t(z_K)=\log a_t(z_K)-\log\widehat z_t.
  \label{eq:p15-log-filter-eval}
\]

For derivatives, store
\[
  \dot{\mathcal H}_{k,ab}
  =
  \dot C_{k,a}\otimes C_{k,b}
  +
  C_{k,a}\otimes \dot C_{k,b},
  \label{eq:p15-Hcoeff-dot}
\]
then build \(\dot Q_k\) and \(\dot E_k(z_k)\) by replacing
\(\mathcal H\) with \(\dot{\mathcal H}\) in \eqref{eq:p15-Qk} and
\eqref{eq:p15-Ek}.  The derivative of the square marginal is
\[
  \dot a_t^{\rm sq}(z_K)
  =
  \sum_{m=1}^D
  P_{m-1}\,\dot T_m\,S_m,
  \label{eq:p15-adot-chain}
\]
where \(\dot T_m=\dot E_m\) for kept coordinates and \(\dot T_m=\dot Q_m\)
for integrated coordinates.  The defensive kept term is fixed, so
\[
  \dot a_t(z_K)=\dot a_t^{\rm sq}(z_K),
  \qquad
  \partial_i\log\widehat p_t(z_K)
  =
  \frac{\dot a_t(z_K)}{a_t(z_K)}
  -
  \frac{\dot{\widehat z}_t}{\widehat z_t}.
  \label{eq:p15-log-filter-dot-eval}
\]
Equations \eqref{eq:p15-H-coeff-shape}--\eqref{eq:p15-log-filter-dot-eval}
are the exact saved-filter representation used by the two-step reference
example.

\section{Saved Branch Data Structures}

The implementation must save the following objects.

\begin{longtable}{p{0.25\linewidth}p{0.68\linewidth}}
\toprule
Object & Required fields \\
\midrule
\texttt{DomainMap} & lower bounds \(\ell_k\), upper bounds \(u_k\), centers
\(a_k\), half-widths \(h_k\), \(\log J\).\\
\texttt{Basis1D} & degree \(p_k\), Legendre recurrence evaluator, mass matrix
\(M_k=I\).\\
\texttt{FitDesign} & Halton points, coordinate primes, uniform weights,
\(N_{\rm fit}\).\\
\texttt{TTCore} & array with shape \((r_{k-1},p_k,r_k)\).\\
\texttt{TTFunction} & ordered cores, basis objects, evaluation routine.\\
\texttt{SquaredTTDensity} & root TT, \(\tau_t\), uniform defensive density,
\(c_t\), \(R_0\), \(\widehat z_t\), square-core environments.\\
\texttt{FilterObject} & numerator representation \(a_t\), normalizer
\(\widehat z_t\), value, log-value, derivative, and log-derivative evaluators.\\
\texttt{BranchObject} & ordering, domain maps, basis degrees, ranks, design
points, weights, ridge, sweep count, initialization, shift, \(\tau_t\),
vectorization order, saved ALS systems.\\
\texttt{SensitivityObject} & derivatives of target values, cores, environments,
normalizers, marginal numerator, and score terms.\\
\bottomrule
\end{longtable}

Saving these objects is not optional.  Without them, an implementation cannot
prove that the derivative pass differentiates the same scalar as the forward
pass.

\section{Forward Filtering Algorithm In Pseudocode}

\begin{enumerate}[label=\textbf{F\arabic*.},leftmargin=*]
\item Choose and save \(B_t\): domain map, Legendre degrees, ranks, Halton
points, weights, ridge, sweep count, initialization rule, \(\epsilon_\tau\),
and \(c_t\).
\item At each fitting point \(z^{(j)}\), compute
\[
  \widetilde q_t(z^{(j)};\alpha)
  =
  \widehat p_{t-1}(\Psi_{t,{\rm old}}(z^{(j)});\alpha)
  f_\alpha(\Psi_{t,x}(z^{(j)})\mid \Psi_{t,{\rm old}}(z^{(j)}))
  g_\alpha(y_t\mid \Psi_{t,x}(z^{(j)}))
  J_t.
\]
\item Form \(y_j=\exp(c_t/2)\sqrt{\widetilde q_t(z^{(j)};\alpha)}\).
\item Run exactly \(S\) bidirectional ridge ALS sweeps using
\eqref{eq:p15-als-row} and \eqref{eq:p15-ridge-normal}.
\item Compute \(R_0\) by \eqref{eq:p15-R-recursion} and
\(\widehat z_t=R_0+\tau_t\).
\item Add \(\log\widehat z_t-c_t\) to the approximate log-likelihood.
\item Build and save the next \texttt{FilterObject} by
\eqref{eq:p15-filter-object}.
\end{enumerate}

\section{Derivative Of Target Values And Previous-Filter Evaluations}

For a positive density value,
\[
  y_j=\exp(c_t/2)\sqrt{\widetilde q_j},
  \qquad
  \dot y_j=\frac12 y_j\,\partial_i\log\widetilde q_j,
  \label{eq:p15-ydot}
\]
because \(c_t\) is frozen.  Here and below a dot denotes
\(\partial_i\).  At time \(t\),
\[
  \partial_i\log\widetilde q_t
  =
  \partial_i\log\widehat p_{t-1}
  +
  \partial_i\log f_\alpha
  +
  \partial_i\log g_\alpha,
  \label{eq:p15-logqdot}
\]
since the branch domain Jacobian is fixed.

The previous filter derivative is evaluated from its saved numerator:
\[
  \widehat p_{t-1}(z_K)=a_{t-1}(z_K)/\widehat z_{t-1},
  \qquad
  \partial_i\log\widehat p_{t-1}
  =
  \frac{\dot a_{t-1}}{a_{t-1}}
  -
  \frac{\dot{\widehat z}_{t-1}}{\widehat z_{t-1}}.
  \label{eq:p15-prev-filter-dot}
\]
This is the operational reason for saving derivative environments for the
next filter object.

In code, the evaluation routine for a new kept-coordinate point performs:
\begin{enumerate}[label=\textbf{E\arabic*.},leftmargin=*]
\item Evaluate the kept-coordinate basis vectors \(b_k(z_k)\).
\item Build \(E_k(z_k)\) and \(\dot E_k(z_k)\) from
\eqref{eq:p15-Ek} and the derivative analogue.
\item Use stored \(Q_k,\dot Q_k\) for integrated coordinates.
\item Form \(T_k,\dot T_k\), prefix products \(P_k\), and suffix products
\(S_k\) by \eqref{eq:p15-prefix-suffix}.
\item Compute \(a_t\), \(\dot a_t\), \(\log\widehat p_t\), and
\(\partial_i\log\widehat p_t\) by
\eqref{eq:p15-log-filter-eval} and \eqref{eq:p15-log-filter-dot-eval}.
\end{enumerate}
No hidden tensor-train operation remains in this evaluator; it is only matrix
products in the squared ranks.

\section{Derivative Through Fixed ALS Sweeps}

Each stored ALS update solves
\[
  N g=d,\qquad N=A^\top W A+\rho I,\qquad d=A^\top W y.
  \label{eq:p15-Ngd}
\]
For the fixed branch, \(W\), \(\rho\), ranks, and fitting points are fixed.
The matrices \(A\) may still depend on \(\alpha\) because their rows contain
environments from other cores, and those cores depend on previous ALS updates.
Differentiating gives
\[
  \dot N=\dot A^\top W A+A^\top W\dot A,
  \qquad
  \dot d=\dot A^\top W y+A^\top W\dot y.
  \label{eq:p15-Nd-dot}
\]

\begin{proposition}[Fixed-rank ridge ALS solve derivative]
If \(N\) is nonsingular, the derivative of the stored update is
\[
  N\dot g=\dot d-\dot N g.
  \label{eq:p15-gdot}
\]
\end{proposition}

\begin{proof}
Differentiate \(Ng=d\).  The product rule gives
\[
  \dot N g+N\dot g=\dot d.
\]
Rearranging gives \eqref{eq:p15-gdot}.  The ridge term helps make \(N\)
nonsingular, but the implementation must still monitor conditioning.
\end{proof}

The derivative of the design row follows from
\[
  A_{j,\cdot}^{(k)}
  =
  (R_j^{(k)})^\top\otimes b_k(z_k^{(j)})^\top\otimes L_j^{(k)}.
\]
The basis values are fixed, so
\[
  \dot A_{j,\cdot}^{(k)}
  =
  (\dot R_j^{(k)})^\top\otimes b_k(z_k^{(j)})^\top\otimes L_j^{(k)}
  +
  (R_j^{(k)})^\top\otimes b_k(z_k^{(j)})^\top\otimes \dot L_j^{(k)}.
  \label{eq:p15-Adot-row}
\]
The environments use ordinary product rules.  For example,
\[
  \dot L_j^{(k)}
  =
  \sum_{m<k}
  G_1\cdots \dot G_m\cdots G_{k-1}.
  \label{eq:p15-Ldot}
\]

\section{Derivative Of Mass Contractions, Marginals, And The Likelihood Scalar}

Differentiate the mass contraction \eqref{eq:p15-R-recursion}.  With fixed
mass matrices,
\[
\begin{aligned}
  \dot R_{k-1}
  =
  \sum_{a,b}M_k[a,b]\Big[
  &\dot C_{k,a}R_kC_{k,b}^{\top}
  +C_{k,a}\dot R_kC_{k,b}^{\top}  \\
  &+C_{k,a}R_k\dot C_{k,b}^{\top}
  \Big],
\end{aligned}
  \label{eq:p15-Rdot}
\]
with \(\dot R_D=0\).  Therefore
\[
  \dot{\widehat z}_t=\dot R_0+\dot\tau_t=\dot R_0.
  \label{eq:p15-zdot}
\]

For the marginal numerator, differentiate the square core:
\[
  \dot H_k(z_k)=\dot G_k(z_k)\otimes G_k(z_k)+G_k(z_k)\otimes \dot G_k(z_k).
  \label{eq:p15-Hdot}
\]
Integrated environments are differentiated by the concrete chain rule
\eqref{eq:p15-adot-chain}.  In array form, the stored matrices have shapes
\[
  Q_k,\dot Q_k,E_k,\dot E_k
  \in
  \R^{r_{k-1}^2\times r_k^2},
  \qquad
  P_k\in\R^{1\times r_k^2},
  \qquad
  S_k\in\R^{r_k^2\times1}.
  \label{eq:p15-env-shapes}
\]
The defensive kept term is fixed, so its derivative is zero in the primary P15
branch.

\begin{proposition}[Normalized approximate filtering]
If \(\phi_t\) is measurable, \(\tau_t>0\), and \(R_0<\infty\), then
\[
  \widehat p_t(z_K)=a_t(z_K)/\widehat z_t
\]
is a nonnegative normalized density on the kept coordinates.
\end{proposition}

\begin{proof}
The density \(\widehat q_t=\phi_t^2+\tau_t2^{-D}\) is nonnegative and has
finite positive integral \(\widehat z_t=R_0+\tau_t\).  Hence
\(\widehat q_t/\widehat z_t\) is a normalized joint density on
\([-1,1]^D\).  The marginal over \(z_K\) is nonnegative and integrates to one
by Tonelli's theorem.
\end{proof}

\begin{proposition}[Same-scalar fixed-branch gradient]
For the saved branch \(B\), the gradient algorithm differentiates exactly
\[
  \widehat\ell_T^{\TT}(\alpha;B)
  =
  \sum_{t=1}^T
  \{\log(R_{t,0}(\alpha;B_t)+\tau_t)-c_t\}.
\]
For one parameter coordinate,
\[
  \partial_i\widehat\ell_T^{\TT}(\alpha;B)
  =
  \sum_{t=1}^T
  \frac{\partial_iR_{t,0}(\alpha;B_t)}
       {R_{t,0}(\alpha;B_t)+\tau_t}.
  \label{eq:p15-final-gradient}
\]
\end{proposition}

\begin{proof}
For the primary branch, \(\dot\tau_t=\dot c_t=0\).  Differentiating the
\(t\)-th summand gives
\[
  \partial_i\{\log(R_{t,0}+\tau_t)-c_t\}
  =
  \frac{\dot R_{t,0}}{R_{t,0}+\tau_t}.
\]
The cores that determine \(R_{t,0}\) are obtained by a fixed finite sequence
of differentiable ridge solves.  Proposition 3 gives the derivative of each
solve, and \eqref{eq:p15-Rdot} gives the derivative of the mass contraction.
The previous-filter dependence is included through
\eqref{eq:p15-prev-filter-dot}.  Therefore the derivative is the derivative of
the same scalar computed by the forward pass.
\end{proof}

\section{Complete Gradient Algorithm In Pseudocode}

\begin{enumerate}[label=\textbf{G\arabic*.},leftmargin=*]
\item Run the forward filter and save every branch object, ALS system, core,
environment, \(R_0\), \(\widehat z_t\), and marginal numerator object.
\item For a parameter coordinate \(\alpha_i\), evaluate
\(\partial_i\log f_\alpha\) and \(\partial_i\log g_\alpha\) at fitting points.
\item At \(t=1\), use the prior derivative.  At \(t>1\), evaluate
\(\partial_i\log\widehat p_{t-1}\) by \eqref{eq:p15-prev-filter-dot}.
\item Form \(\dot y_j\) by \eqref{eq:p15-ydot}.
\item Replay the fixed ALS sweeps.  At every core update, form
\(\dot A\), \(\dot N\), \(\dot d\), and solve \eqref{eq:p15-gdot}.
\item Compute \(\dot R_0\) by \eqref{eq:p15-Rdot}.
\item Add \(\dot R_0/(R_0+\tau_t)\) to the score.
\item Build derivative environments for the next \texttt{FilterObject}: save
\(\mathcal H_{k,ab}\), \(\dot{\mathcal H}_{k,ab}\), \(Q_k\), \(\dot Q_k\),
\(\widehat z_t\), \(\dot{\widehat z}_t\), and \(\tau_t2^{-|K|}\).  Its
evaluation routine is exactly \textbf{E1}--\textbf{E5}.
\end{enumerate}

This recipe is intentionally replay based.  It is longer than simply applying
automatic differentiation, but every derivative is tied to a stored object in
the forward pass.

\section{Numerical Stabilization And Failure Diagnostics}

The implementation must report:
\[
  \frac{\|A g-y\|_2}{\|y\|_2},\qquad
  \max_k r_k,\qquad
  \max_{\text{ALS updates}}\kappa(N_k),
  \qquad
  \frac{\tau_t}{R_0+\tau_t}.
  \label{eq:p15-diagnostics}
\]
It must also report whether any fitted or marginal numerator evaluation is
nonfinite.  A large residual means the square-root density was not fitted
well.  A large condition number means the fixed linear algebra is unstable.
An overly large defensive-mass ratio means the uniform floor is dominating the
filter.  These diagnostics can veto a run, but they do not prove posterior
accuracy when they look good.

\section{Finite-Difference Parity Test}

For the saved branch \(B_0\), test the scalar
\[
  \widehat\ell_T^{\TT}(\alpha;B_0)
  =
  \sum_t\{\log(R_{t,0}(\alpha;B_0)+\tau_t)-c_t\}.
  \label{eq:p15-parity-scalar}
\]
For each tested coordinate \(i\), compute
\[
  D_i(h)=
  \frac{
  \widehat\ell_T^{\TT}(\alpha_0+h e_i;B_0)
  -
  \widehat\ell_T^{\TT}(\alpha_0-h e_i;B_0)}
  {2h},
  \label{eq:p15-centered-difference}
\]
using
\[
  h\in\{10^{-3},10^{-4},10^{-5},10^{-6}\}\max(1,|\alpha_{0,i}|).
  \label{eq:p15-h-ladder}
\]
Report
\[
  e_i(h)=
  \frac{|D_i(h)-\partial_i\widehat\ell_T^{\TT}(\alpha_0;B_0)|}
       {1+|D_i(h)|+|\partial_i\widehat\ell_T^{\TT}(\alpha_0;B_0)|}.
  \label{eq:p15-parity-error}
\]
The branch objects, fitting points, ranks, shift, defensive mass, and sweep
count must remain unchanged at \(\alpha_0\pm h e_i\).  If the smallest
\(e_i(h)\) is larger than \(10^{-4}\) in the minimal example, the reference
implementation fails the same-scalar test.

\section{Minimal Runnable Example}

The P15 reference script implements a two-step scalar nonlinear model:
\[
  x_0\sim N(0,\sigma_0^2),\qquad
  x_1\mid x_0\sim N(\alpha x_0,\sigma_\eta^2),
  \label{eq:p15-toy-transition}
\]
\[
  y_1\mid x_1\sim N(x_1^2,\sigma_\epsilon^2).
  \label{eq:p15-toy-observation}
\]
The second step uses the saved approximate filter from the first step:
\[
  \widehat p_1(x_1;\alpha)=a_1(x_1;\alpha)/\widehat z_1(\alpha),
  \qquad
  x_2\mid x_1\sim N(\alpha x_1,\sigma_\eta^2),
  \qquad
  y_2\mid x_2\sim N(x_2^2,\sigma_\epsilon^2).
  \label{eq:p15-toy-step2}
\]
Thus the example exercises the recursive carry-forward object, not only the
one-step normalizer.

At each step the augmented variable is \(r=(x_t,x_{t-1})\), so \(D=2\).  The
script uses \([-3,3]^2\), degree \(5\), \(N_{\rm fit}=96\), fixed rank one,
ridge \(10^{-8}\), five bidirectional sweeps, and
\(\epsilon_\tau=10^{-10}\).  The parameter is \(\alpha\).

The derivative of the transformed log density is
\[
  \partial_\alpha\log\widetilde q(z;\alpha)
  =
  \frac{(x_1-\alpha x_0)x_0}{\sigma_\eta^2},
  \label{eq:p15-toy-logqdot}
\]
because the prior, likelihood, and fixed Jacobian do not depend on
\(\alpha\).  The script then applies \eqref{eq:p15-ydot},
\eqref{eq:p15-gdot}, and \eqref{eq:p15-final-gradient}.

At the second step, the derivative of the transformed log density includes the
saved-filter term:
\[
  \partial_\alpha\log\widetilde q_2
  =
  \partial_\alpha\log\widehat p_1(x_1;\alpha)
  +
  \frac{(x_2-\alpha x_1)x_1}{\sigma_\eta^2}.
  \label{eq:p15-toy-step2-dot}
\]
The script obtains \(\partial_\alpha\log\widehat p_1\) from
\eqref{eq:p15-log-filter-dot-eval}.  This is the concrete recursive path that
the one-step P14 note did not test.

The run produced:
\[
\begin{array}{ll}
c_1&=-0.445460193510,\\
c_2&=-1.496961434280,\\
\widehat\ell&=-0.569442222488,\\
\partial_\alpha\widehat\ell&=-1.004731548060,\\
\text{step 1 fit residual}&=3.905823\times10^{-1},\\
\text{step 2 fit residual}&=3.544318\times10^{-1},\\
\max\kappa(N)\text{ at step 1}&=2.977710,\\
\max\kappa(N)\text{ at step 2}&=2.914195.
\end{array}
\]
The centered-difference relative errors were
\[
\begin{array}{ccc}
h & D(h) & e(h)\\
10^{-3} & -1.004731310930 & 7.879523\times10^{-8}\\
10^{-4} & -1.004731545700 & 7.856375\times10^{-10}\\
10^{-5} & -1.004731548040 & 7.604872\times10^{-12}\\
10^{-6} & -1.004731548630 & 1.879178\times10^{-10}.
\end{array}
\]
The parity check passes.  The large-ish fit residual is not hidden: this toy
example is a same-scalar gradient and implementability check, not a claim that
rank one is an accurate filter for the nonlinear density.

\section{Relation To Zhao--Cui And What Remains To Be Tested}

Zhao--Cui use sequential tensor-train approximations of the joint density over
\((x_t,\theta,x_{t-1})\), squared density representations of the form
\(\phi^2+\tau\lambda\), normalizers, marginalization, and conditional
transport maps.  P15 keeps that filtering architecture but freezes the
numerical branch and replaces the adaptive approximation machinery by one
explicit fixed-rank ridge ALS path so that a same-scalar derivative can be
written and tested.

What this note proves is limited but useful:
\[
  \text{fixed branch}+\text{squared density}+\text{finite positive mass}
  \Longrightarrow
  \text{normalized approximate filter},
\]
and
\[
  \text{fixed branch}+\text{differentiable stored solves}
  \Longrightarrow
  \nabla_\alpha\widehat\ell_T^{\TT}(\alpha;B)
  \text{ for the same scalar.}
\]
What remains empirical is accuracy and efficiency on a target high-dimensional
model: fit residuals, rank growth, conditioning, defensive-mass ratio,
finite-difference parity, and comparison against fixed sparse-grid Gaussian
projection filtering.

\appendix
\section{Source And Code Anchors}

The main text is self-contained.  These anchors explain provenance:
\begin{itemize}[leftmargin=*]
\item Zhao--Cui state-space model equations (1)--(3): source for the SSM
notation.
\item Zhao--Cui equations (9)--(12) and Algorithm 1: source for the sequential
joint density over \((x_t,\theta,x_{t-1})\) and the integration architecture.
\item Zhao--Cui equation (13), Lemma 1, Algorithm 2, and Proposition 2: source
for squared-TT density and marginal construction.
\item Zhao--Cui Section 4.1: source for the evidence/normalizer role.
\item Cui--Dolgov Sections 2--3: source context for squared inverse
Rosenblatt and functional tensor-train transport.
\item Inspected code paths: \texttt{models/full\_sol.m} updates
\(\log(\texttt{sirt.z})-\texttt{const}\), and
\texttt{@TTSIRT/marginalise.m} computes
\(\texttt{obj.z}=\texttt{obj.fun\_z}+\texttt{obj.tau}\).
\end{itemize}

The exact Legendre basis, Halton design, fixed-rank ridge ALS protocol, frozen
shift rule, and reference Python example are P15 design choices.  They are not
claimed to be Zhao--Cui's adaptive MATLAB implementation.

\section{Symbol Table And Array Shapes}

\begin{longtable}{p{0.23\linewidth}p{0.67\linewidth}}
\toprule
Symbol & Meaning \\
\midrule
\(D\) & augmented dimension of \(r_t=(x_t,\vartheta,x_{t-1})\).\\
\(z\) & reference coordinate in \([-1,1]^D\).\\
\(\Psi_t\) & affine map from reference coordinate to physical coordinate.\\
\(b_{k,a}\) & normalized Legendre basis function for coordinate \(k\).\\
\(C_{k,a}\) & coefficient slice, shape \(r_{k-1}\times r_k\).\\
\(G_k(z_k)\) & TT factor, shape \(r_{k-1}\times r_k\).\\
\(\phi_t\) & square-root TT function.\\
\(H_k\) & square core \(G_k\otimes G_k\).\\
\(R_k\) & suffix mass-contraction environment.\\
\(A^{(k)}\) & design matrix for the ALS update of core \(k\).\\
\(N_k\) & ridge normal matrix \(A^\top W A+\rho I\).\\
\(a_t\) & marginal numerator for the next filter.\\
\(\widehat z_t\) & approximate one-step normalizer.\\
\(B_t\) & all frozen branch objects at time \(t\).\\
\bottomrule
\end{longtable}

\end{document}
